%==============================================================================
% Dàn ý chi tiết - Chapter 3: Phương pháp đề xuất
%==============================================================================

\chapter{Phương pháp đề xuất}

%-----------------------------------------------------------------------------
\section{Tổng quan hệ thống}
%-----------------------------------------------------------------------------
\begin{itemize}
    \item Sơ đồ kiến trúc tổng thể của hệ thống
    \begin{itemize}
        \item Các thành phần chính
        \item Luồng dữ liệu (data flow)
    \end{itemize}
    \item Mô tả ngắn gọn các bước trong pipeline
    \begin{enumerate}
        \item Tiền xử lý dữ liệu đầu vào
        \item Sinh câu hỏi (Single-Prompt hoặc Multi-Agent)
        \item Tạo hình ảnh và OCR (evaluation)
        \item Đánh giá chất lượng
    \end{enumerate}
    \item Các đầu vào và đầu ra của hệ thống
    \item So sánh với các phương pháp hiện có
\end{itemize}

%-----------------------------------------------------------------------------
\section{Bộ dữ liệu}
%-----------------------------------------------------------------------------

\subsection{Tập dữ liệu RACE}
\begin{itemize}
    \item Giới thiệu về RACE dataset:
    \begin{itemize}
        \item Nguồn gốc và mục đích
        \item Cấu trúc dữ liệu (article, questions, options, answer)
        \item Số lượng mẫu (train/dev/test)
        \item Đặc điểm: Reading comprehension từ kỳ thi Trung Quốc
    \end{itemize}
    \item Lý do chọn RACE:
    \begin{itemize}
        \item Đa dạng chủ đề
        \item Độ khó phân hóa tốt
        \item Phổ biến trong nghiên cứu
    \end{itemize}
    \item Thống kê dữ liệu:
    \begin{itemize}
        \item Số lượng bài đọc
        \item Số lượng câu hỏi
        \item Độ dài trung bình
    \end{itemize}
\end{itemize}

\subsection{Tiền xử lý và chuẩn hóa dữ liệu}
\begin{itemize}
    \item Các bước tiền xử lý:
    \begin{itemize}
        \item Loại bỏ noise
        \item Chuẩn hóa encoding
        \item Xử lý special characters
    \end{itemize}
    \item Tách văn bản thành đơn vị xử lý
    \item Format output mong muốn
    \item Data augmentation (nếu có)
\end{itemize}

%-----------------------------------------------------------------------------
\section{Phương pháp sinh câu hỏi trắc nghiệm}
%-----------------------------------------------------------------------------

\subsection{Phương pháp đơn prompt (Single-Prompt)}
\begin{itemize}
    \item Mô tả phương pháp:
    \begin{itemize}
        \item Sử dụng một prompt duy nhất để sinh câu hỏi cho tất cả cấp độ
    \end{itemize}
    \item Prompt template:
    \begin{itemize}
        \item System prompt
        \item User prompt với context
    \end{itemize}
    \item Các tham số:
    \begin{itemize}
        \item Temperature
        \item Max tokens
        \item N
    \end{itemize}
    \item Ưu điểm:
    \begin{itemize}
        \item Đơn giản, dễ implement
        \item Chi phí thấp
    \end{itemize}
    \item Nhược điểm:
    \begin{itemize}
        \item Khó kiểm soát chất lượng
        \item Thiếu tính chuyên môn hóa
    \end{itemize}
\end{itemize}

\subsection{Phương pháp đa tác tử (Multi-Agent)}
\begin{itemize}

\subsection{Tác tử phân tích (Analyzer Agent)}
\begin{itemize}
    \item Nhiệm vụ:
    \begin{itemize}
        \item Phân tích văn bản đầu vào
        \item Xác định các thực thể, sự kiện quan trọng
        \item Trích xuất thông tin key
    \end{itemize}
    \item Prompt design:
    \begin{itemize}
        \item Role và instructions
        \item Output format
    \end{itemize}
    \item Ví dụ output của Analyzer
\end{itemize}

\subsection{Tác tử sinh câu hỏi theo cấp độ (Generator Agents)}
\begin{itemize}
    \item Generator Level 1 (Retrieval):
    \begin{itemize}
        \item Nhiệm vụ: Sinh câu hỏi yêu cầu nhớ lại thông tin
        \item Đặc điểm: Câu hỏi có thể trả lời trực tiếp từ văn bản
    \end{itemize}
    \item Generator Level 2 (Inference \& Synthesis):
    \begin{itemize}
        \item Nhiệm vụ: Sinh câu hỏi yêu cầu suy luận
        \item Đặc điểm: Cần tổng hợp thông tin từ nhiều phần
    \end{itemize}
    \item Generator Level 3 (Critical Reasoning):
    \begin{itemize}
        \item Nhiệm vụ: Sinh câu hỏi yêu cầu đánh giá, phản biện
        \item Đặc điểm: Cần kiến thức ngoài văn bản hoặc quan điểm cá nhân
    \end{itemize}
    \item Mỗi generator có prompt riêng biệt
    \item Parallel generation cho efficiency
\end{itemize}

\subsection{Tác tử tổng hợp (Merge Agent)}
\begin{itemize}
    \item Nhiệm vụ:
    \begin{itemize}
        \item Tổng hợp kết quả từ các generator
        \item Loại bỏ trùng lặp
        \item Chuẩn hóa format đầu ra
    \end{itemize}
    \item Prompt design
    \item Xử lý conflicts
\end{itemize}

\subsection{Quy trình hoạt động (Workflow)}
\begin{itemize}
    \item Sơ đồ luồng xử lý (flowchart)
    \begin{itemize}
        \item Bước 1: Analyzer xử lý văn bản
        \item Bước 2: Gửi cho 3 generator (song song)
        \item Bước 3: Merge agent tổng hợp
        \item Bước 4: Output cuối cùng
    \end{itemize}
    \item State management trong LangGraph
    \item Error handling
    \item Retry mechanism
\end{itemize}
\end{itemize}

%-----------------------------------------------------------------------------
\section{Thiết kế prompt theo cấp độ Bloom}
%-----------------------------------------------------------------------------

\subsection{Prompt cấp độ 1 -- Truy xuất}
\begin{itemize}
    \item System prompt:
    \begin{itemize}
        \item Định nghĩa role
        \item Các nguyên tắc sinh câu hỏi Retrieval
    \end{itemize}
    \item User prompt template:
    \begin{itemize}
        \item Định dạng input
        \item Yêu cầu cụ thể
    \end{itemize}
    \item Few-shot examples:
    \begin{itemize}
        \item Ví dụ văn bản $\rightarrow$ câu hỏi
    \end{itemize}
    \item Output format specification
    \item Constraints và注意事项
\end{itemize}

\subsection{Prompt cấp độ 2 -- Suy luận và tổng hợp}
\begin{itemize}
    \item System prompt:
    \begin{itemize}
        \item Định nghĩa role
        \item Các nguyên tắc sinh câu hỏi Inference/Synthesis
    \end{itemize}
    \item User prompt template
    \item Few-shot examples
    \item Output format specification
    \item Constraints và lưu ý
\end{itemize}

\subsection{Prompt cấp độ 3 -- Lập luận phản biện}
\begin{itemize}
    \item System prompt:
    \begin{itemize}
        \item Định nghĩa role
        \item Các nguyên tắc sinh câu hỏi Critical Reasoning
    \end{itemize}
    \item User prompt template
    \item Few-shot examples
    \item Output format specification
    \item Constraints và lưu ý
    \item Cách xử lý câu hỏi mở (cần có evidence)
\end{itemize}

%-----------------------------------------------------------------------------
\section{Quy trình OCR}
%-----------------------------------------------------------------------------

\subsection{Sinh ảnh sạch từ văn bản (Clean Image Generation)}
\begin{itemize}
    \item Mục đích: Tạo baseline images cho OCR
    \item Công cụ sử dụng:
    \begin{itemize}
        \item Pillow/PIL
        \item ReportLab
    \end{itemize}
    \item Quy trình:
    \begin{itemize}
        \item Convert text to image
        \item Font selection
        \item Layout configuration
    \end{itemize}
    \item Tham số:
    \begin{itemize}
        \item Resolution (DPI)
        \item Font size
        \item Image format (PNG, JPEG)
    \end{itemize}
\end{itemize}

\subsection{Tăng cường ảnh (Image Augmentation với Augraphy)}
\begin{itemize}
    \item Giới thiệu Augraphy:
    \begin{itemize}
        \item Thư viện cho document augmentation
    \end{itemize}
    \item Các kỹ thuật augmentation:
    \begin{itemize}
        \item Noise addition (Gaussian, Poisson)
        \item Blur effects
        \item Rotation và skew
        \item Paper texture
        \item Ink bleeding
        \item Folding effects
    \end{itemize}
    \item Pipeline augmentation:
    \begin{itemize}
        \item Sequential augmentation steps
        \item Randomization parameters
    \end{itemize}
    \item Mục đích: Mô phỏng điều kiện scanning thực tế
\end{itemize}

\subsection{Nhận dạng văn bản từ ảnh (Text Recognition)}
\begin{itemize}
    \item Phương pháp Traditional OCR:
    \begin{itemize}
        \item EasyOCR
        \item PaddleOCR
    \end{itemize}
    \item Phương pháp VLM-based:
    \begin{itemize}
        \item GPT-4V
        \item Claude Vision
    \end{itemize}
    \item Pipeline xử lý:
    \begin{itemize}
        \item Preprocessing
        \item Recognition
        \item Post-processing
    \end{itemize}
    \item So sánh các phương pháp
\end{itemize}

%-----------------------------------------------------------------------------
\section{Khung đánh giá (Evaluation Framework)}
%-----------------------------------------------------------------------------

\subsection{Đánh giá vi mô (Micro Evaluation)}
\begin{itemize}

\subsection{Tính giải được (Solvability)}
\begin{itemize}
    \item Định nghĩa: Câu hỏi có thể trả lời được từ văn bản
    \item Phương pháp đánh giá:
    \begin{itemize}
        \item Sử dụng LLM để kiểm tra answerability
        \item Prompt: "Có thể trả lời câu hỏi này từ văn bản không?"
    \end{itemize}
    \item Metric: Percentage of solvable questions
\end{itemize}

\subsection{Chất lượng phương án nhiễu (Distractor Quality)}
\begin{itemize}
    \item Định nghĩa: Các phương án sai phải hợp lý
    \item Phương pháp đánh giá:
    \begin{itemize}
        \item Kiểm tra plausibility của distractors
        \item Không quá dễ loại trừ
        \item Liên quan đến chủ đề
    \end{itemize}
    \item Sử dụng LLM để đánh giá
    \item Rubric đánh giá
\end{itemize}

\subsection{Tính phù hợp cấp độ (Alignment)}
\begin{itemize}
    \item Định nghĩa: Câu hỏi đúng cấp độ Bloom
    \item Phương pháp đánh giá:
    \begin{itemize}
        \item Sử dụng LLM để classify cấp độ
        \item So sánh với cấp độ mong muốn
    \end{itemize}
    \item Metric: Alignment score
\end{itemize}
\end{itemize}

\subsection{Đánh giá vĩ mô (Macro Evaluation)}
\begin{itemize}
    \item Định nghĩa: Đánh giá tổng thể chất lượng câu hỏi
    \item Các tiêu chí:
    \begin{itemize}
        \item Grammar and fluency
        \item Clarity
        \item Difficulty appropriateness
        \item Overall quality
    \end{itemize}
    \item Phương pháp:
    \begin{itemize}
        \item LLM-as-a-Judge với detailed rubric
        \item Pairwise comparison (nếu có)
    \end{itemize}
    \item Thang điểm 1-5
    \item Statistical analysis
\end{itemize}

\subsection{Đánh giá định dạng đầu ra (Format Evaluation)}
\begin{itemize}
    \item Kiểm tra format chuẩn:
    \begin{itemize}
        \item Question text
        \item 4 options (A, B, C, D)
        \item Correct answer
    \end{itemize}
    \item Tỷ lệ lỗi:
    \begin{itemize}
        \item Syntax error rate
        \item Missing fields
    \end{itemize}
    \item Over-generation và Under-generation:
    \begin{itemize}
        \item Số câu hỏi sinh ra so với yêu cầu
    \end{itemize}
    \item Automated parsing và validation
\end{itemize}

\subsection{Đánh giá OCR (CER và WER)}
\begin{itemize}
    \item Character Error Rate (CER):
    \begin{itemize}
        \item Công thức tính
        \item Ý nghĩa
    \end{itemize}
    \item Word Error Rate (WER):
    \begin{itemize}
        \item Công thức tính
        \item Ý nghĩa
    \end{itemize}
    \item So sánh CER/WER giữa:
    \begin{itemize}
        \item Ảnh sạch vs ảnh augmentation
        \item Traditional OCR vs VLM OCR
    \end{itemize}
    \item Statistical significance
\end{itemize}

%-----------------------------------------------------------------------------
\section{Ứng dụng minh họa -- Openlet}
%-----------------------------------------------------------------------------

\subsection{Kiến trúc hệ thống}
\begin{itemize}
    \item Frontend:
    \begin{itemize}
        \item Next.js
        \item UI components
    \end{itemize}
    \item Backend:
    \begin{itemize}
        \item Firebase (Firestore, Storage)
        \item API routes
    \end{itemize}
    \item AI Integration:
    \begin{itemize}
        \item OpenAI API
        \item Claude API
    \end{itemize}
    \item Database schema
    \item Authentication
\end{itemize}

\subsection{Luồng hoạt động}
\begin{itemize}
    \item User workflow:
    \begin{enumerate}
        \item User upload text/document
        \item Select difficulty level(s)
        \item System generates questions
        \item User can edit/review
        \item Export to PDF/Quiz format
    \end{enumerate}
    \item API endpoints
    \item Error handling
    \item User feedback loop
\end{itemize}
\end{itemize}
